% Chapter 1: Introduction

\section{Background}

Botswana's agricultural sector, particularly livestock farming, plays a crucial role in the country's economy and food security. However, traditional livestock management faces significant challenges including limited real-time monitoring capabilities, disease outbreak detection delays, and inefficient resource allocation. The integration of Internet of Things (IoT) technologies and intelligent systems offers promising solutions to these challenges.

\section{Problem Statement}

Current livestock management systems in Botswana suffer from several critical limitations:

\begin{itemize}
    \item Lack of real-time health monitoring for individual animals
    \item Delayed disease detection leading to economic losses
    \item Inefficient data governance and privacy concerns
    \item Limited decision support for farmers and veterinarians
    \item Poor integration between data collection and analysis systems
\end{itemize}

These challenges result in reduced productivity, increased mortality rates, and suboptimal resource utilization in the agricultural sector.

\section{Motivation}

The motivation for this project stems from the urgent need to modernize livestock management practices in Botswana through:

\begin{itemize}
    \item Leveraging emerging IoT and AI technologies
    \item Addressing data governance and privacy concerns
    \item Providing actionable insights for stakeholders
    \item Creating a scalable and sustainable solution
\end{itemize}

\section{Aim}

The primary aim of this project is to develop an autonomous telemetry and governance framework for livestock management in Botswana that integrates real-time data collection, intelligent analysis, and robust data governance mechanisms.

\section{Objectives}

The specific objectives of this project are:

\begin{enumerate}
    \item To design and implement a Third Normal Form (3NF) database schema for livestock telemetry data
    \item To develop a multi-agent orchestration system for autonomous diagnostic routing
    \item To implement role-based access control mechanisms for privacy protection
    \item To create a user-friendly web interface for data visualization and management
    \item To validate the system through comprehensive testing and evaluation
\end{enumerate}

\section{Scope}

The project scope includes:

\subsection{In Scope}
\begin{itemize}
    \item Development of a 3NF validation engine for data integrity
    \item Implementation of multi-agent graph orchestrator
    \item Creation of a React-based web frontend
    \item Integration of SQLite database for edge deployment
    \item Implementation of RBAC for privacy protection
    \item Testing and validation of all components
\end{itemize}

\subsection{Out of Scope}
\begin{itemize}
    \item Hardware sensor development
    \item Mobile application development (deferred to future work)
    \item Integration with external agricultural databases
    \item Machine learning model training (deferred to future work)
    \item Blockchain implementation (deferred to future work)
\end{itemize}

\section{Technologies}

The project utilizes the following technologies:

\begin{table}[H]
\centering
\begin{tabular}{|l|l|l|}
\hline
\textbf{Component} & \textbf{Technology} & \textbf{Purpose} \\
\hline
Backend & Python 3.10+ & Core application logic \\
Web Framework & Flask 3.0.0 & REST API and web server \\
Database & SQLite & Edge-gateway persistence \\
Frontend & React 18+ & Modern web UI \\
Build Tool & Vite & Fast frontend development \\
Styling & Tailwind CSS & Utility-first CSS framework \\
\hline
\end{tabular}
\caption{Technology Stack}
\end{table}

\section{Risks}

\subsection{Technical Risks}
\begin{itemize}
    \item Database performance issues with large datasets
    \item Network connectivity challenges in rural areas
    \item Integration complexity between components
\end{itemize}

\subsection{Project Risks}
\begin{itemize}
    \item Time constraints for comprehensive testing
    \item Limited access to real-world testing environments
    \item Potential scope creep during development
\end{itemize}

\section{Structure of Dissertation}

This dissertation is organized as follows:

\begin{itemize}
    \item \textbf{Chapter 2: Literature Review} - Examines existing research and systems
    \item \textbf{Chapter 3: Methodology} - Describes the research approach
    \item \textbf{Chapter 4: Analysis} - Presents system requirements and analysis
    \item \textbf{Chapter 5: Design} - Details system architecture and design
    \item \textbf{Chapter 6: Implementation} - Describes the implementation process
    \item \textbf{Chapter 7: Testing} - Presents testing methodology and results
    \item \textbf{Chapter 8: Results} - Discusses system performance and outcomes
    \item \textbf{Chapter 9: Discussion} - Analyzes findings and implications
    \item \textbf{Chapter 10: Conclusion} - Summarizes contributions and future work
\end{itemize}
